\documentclass[11pt]{article}
\usepackage{graphicx} % Required for inserting images
\usepackage[utf8]{inputenc}
\usepackage{amsmath,amssymb}
\usepackage{amsthm}
\usepackage{hyperref}
\usepackage{tikz-cd}
\usetikzlibrary{arrows}
\usepackage{geometry}
\usepackage{verbatim}
\geometry{margin=1in}
\theoremstyle{plain}
\newtheorem{theorem}{Theorem}
\newtheorem{proposition}{Proposition}
\newtheorem{lemma}{Lemma}
\newtheorem{corollary}{Corollary}
\theoremstyle{definition}
\newtheorem{definition}{Definition}
\theoremstyle{remark}
\newtheorem{example}{Example}
\newtheorem*{remark}{Remark}
\usepackage{stmaryrd} 
\usepackage{natbib}            
\usepackage{amsfonts}
\newtheorem{result}{Result}

\title{Flower Placement Optimization Problem}
\author{Gordon Tin}
\date{April 2025}

\begin{document}

\maketitle

\section{Problem Statement}

Given an $n \times m$ king-graph $G$, define $A$ and $B$ as two disjoint sets of vertices, representing two different placements within $G'$, and let the function $p(x,y)$ be defined as: 
\[p(x,y)= \frac{0.3xy}{0.5(x+y)^{2}-0.5(x+y)}\]
where $x$ and $y$ represent the number of adjacent vertices in $A$ and $B$, respectively. Let $G'$ be the induced subgraph of $G$ after the removal of $k$ vertices from $G$ at random and, following this, the removal of all connected components of size $\leq 2$.
Our objective is to determine the optimal placement of vertices $A$ and $B$ such that the total probability across all vertices in $G'$ is maximized: \[\sum_{x,y \in G'} p(x,y)\] 
As \(p(x,y)\) declines, the diminishing returns effect in placement incentivizes optimal placement, rather than naively increasing x and y indiscriminately.

We will first reduce the problem to an instance of CircuitSAT to show that it is, in fact, NP. Next, we will show that it is NP-Complete. 

\section{Proof of NP}

To reduce the optimization problem to a decision problem, we want to find a choice of disjoint sets $A, b \subset V(G')$ such that the total payoff $\sum_{v \in G'} p(x_v, y_v) \geq T$, where $T$ is the target value. Solving a sequence of such decision problems, we can identify the maximum achievable sum. 

Let $C$ be a Boolean circuit that outputs $1$ (true) if and only if there exists a valid assignment, i.e., a choice of $A$ and $B$, with a total payoff $\geq T$. Let also $a_i$ be a Boolean variable true exactly when vertex $i$ is in $A$, and $b_i$ be a Boolean variable true exactly when vertex $i$ is in $B$. These variables are the inputs to our circuit, with those vertices in the complement of $G'$ being set to false as they cannot be meaningfully chosen.

We introduce a constraint of disjointedness, where no vertex can belong to $A \cap B$; formally, for each vertex $i$, $\neg (a_i \land b_i)$.

For each vertex $v \in G'$, let $N(v)$ denote the set of neighbors of $v$ in a graph $G'$. Let:
\begin{enumerate}
    \item $x_v = \sum_{u \in N(v)} a_u$
    \item $y_v = \sum_{u \in N(v)} b_u$
\end{enumerate}

These sums are bounded (at most 8 for an interior vertex's degree), so we can represent each of $x_v$ and $y_v$ with 4 bits (to cover 0-8). The circuit will enforce that these encoded values indeed equal the sum of the corresponding neighbor-selection bits by introducing auxiliary Boolean variables to encode the binary digits of $x_v$ and add constraints to link them to the $a_u, b_u$ of neighbors, representable as a subcircuit that outputs the binary sum of inputs $a_u$ for $u \in N(v)$ and similarly for $b_u$.

Next, we need to incorporate the diminishing return function $p(x,y)$ for each vertex into our logic. Note that $p(x,y)$ is a rational function, so, to encode this, we can clear the denominators by multiplying both sides of the equation
\[z_v= \frac{0.3x_v y_v}{0.5(x_v+y_v)^{2}-0.5(x_v+y_v)}\]
by the denominator and then scaling by 10:
\[5((x_v + y_v)^2 - (x_v + y_v)) \cdot z_v = 3 \cdot (x_v \cdot y_v)\]
This gives us a polynomial constraint that forces $z_v$ to equal the value of the diminishing returns function given $x_v$ and $y_v$, which is easier to encode in a Boolean logic setting. As $x_v, y_v$ are small, the relationship can be enforced by a small subcircuit (with multipliers and adders for the arithmetic). Importantly, as $x_v + y_v$ grows, the ratio $z_v$ grows slower, encapsulating the property of diminishing returns.

Finally, we need to enforce that the sum of all vertex payoffs is at least $T$. Let $S = \sum_{v \in G'} z_v$ be the total payoff. We compute this sum $S$ from all of the individual $z_v$ values with pairwise additions since there are a polynomial number of terms. Then we assert that $S \geq T$. In a Boolean circuit, a comparator subcircuit can output 1 if $S$'s binary value is $\geq$ the binary representation of $T$. We then constrain the circuit's output to be true if and only if this comparator outputs 1 and all of the earlier constraints hold true. By the NP-completeness of CircuitSAT, this shows that our decision version is in NP.

Now that we have proved that this problem is in NP, we will now prove that the problem is NP-Complete.

\section{Monotone 2-in-4 SAT}

To prove that our problem is NP-Complete, we reduce from monotone 2-in-4 SAT, which is known to be NP-Complete \citep{KaraKratochvilWood2007}. The problem is defined in \citep{scheffler2022distance} as follows: given a boolean formula $\phi$ in conjunctive normal form, where every clause contains four nonnegated variables, is there an assignment of the variables of $\phi$ such that each clause of $\phi$ contains two variables with value 1 and two variables with value 0?

\subsection{Path Lemmas}

We begin by investigating simple paths. We provide some simple lemmas regarding paths.
\begin{lemma} \label{lem:pathlength}
    For a path of length $k$, there are a maximum of $\lfloor\frac{k-1}{2}\rfloor$ contributing vertices. 
\end{lemma}
\begin{proof}
    Trivial.
\end{proof}
Now, for our gadget, we work with a four-armed gadget, each arm being of length $4k+3$, including the center $c_j$. 
\[
\begin{tikzcd}[
  row sep={2.5em,between origins},
  column sep={2.5em,between origins}
]
x_{i_1} \arrow[dr, dash]
&&&&&&&& x_{i_2} \arrow[dl, dash] \\
& z_{1,4k+1} \arrow[dr, dash]
&&&&&& z_{2,4k+1} \arrow[dl, dash] \\
&& \cdots \arrow[dr, dash]
&&&& \cdots \arrow[dl, dash] \\
&&& z_{1,1} \arrow[dr, dash]
&& z_{2,1} \arrow[dl, dash] \\
&&&& c_j \arrow[dl, dash] \arrow[dr, dash] \\
&&& z_{4,1} \arrow[dl, dash]
&& z_{3,1} \arrow[dr, dash] \\
&& \cdots \arrow[dl, dash]
&&&& \cdots \arrow[dr, dash] \\
& z_{4,4k+1} \arrow[dl, dash]
&&&&&& z_{3,4k+1} \arrow[dr, dash] \\
x_{i_4} &&&&&&&& x_{i_3}
\end{tikzcd}
\]
By Lemma \ref{lem:pathlength}, each arm should be able to provide individually $2k+1$ vertices

\begin{lemma}
    
\end{lemma}

\bibliographystyle{plainnat}
\bibliography{refs}

\end{document}

